\section{Explanation-closure audit}
\label{app:explanation-audit}

After support reduction, normalization handles the noncommuting support-two shapes and leaves configurations containing commuting support-two blocks. The successful final step treats each remaining block's assignment of its two targets to its two frames as an independent feasible permutation. For each of 378 irreducible geometries, the complete target-letter domain at the two active qubits contains $2^{24}=16{,}777{,}216$ states. Every state admits an alternative with cost change $\Delta\leq0$ and strictly fewer commuting support-two blocks. Induction on the count of such blocks reduces every optimum to the union of the weight-one arbitrary-anchor, enlarged-borrow, and hybrid families. This proves Eq.~\eqref{eq:classification}.

The following attempted routes remain relevant because their failure prevents a misleadingly short proof.

\begin{table}[h]
\centering
\caption{Retained negative proof evidence.}
\begin{tabularx}{\textwidth}{>{\raggedright\arraybackslash}p{.27\textwidth}>{\raggedright\arraybackslash}X>{\raggedright\arraybackslash}p{.22\textwidth}}
\toprule
Attempt & Exact outcome & Consequence\\
\midrule
Replacement-orbit descent & Ten candidate images re-entered the same qubit pair; the best available changes tied at zero and no strict global descent was obtained. & The orbit argument does not close the sector.\\
Single-pinner normalization & Confirmed over 133,349,376 composite configurations, but composition into a global chain was not proved. & A valid local lemma is not a global theorem.\\
Two-coordinate chain representation & Refuted by an exact three-tag configuration with multiple irreducible commuting support-two blocks; no new minimizing family was found. & That proof premise fails, while the per-block-permutation theorem remains intact.\\
\bottomrule
\end{tabularx}
\end{table}

These outcomes are compatible: a local normalization can be correct, its proposed chain representation can fail, and a different complete-domain induction can still prove the classification. No failed route is counted as support for the successful theorem.
